\subsection{Phân tích các bên liên quan}

Việc phát triển và triển khai Hệ thống Thương mại Điện tử cho Bán lẻ Điện thoại tích hợp Trợ lý Ảo AI Agent là một bước tiến công nghệ quan trọng, mang lại những lợi ích thiết thực và giá trị gia tăng rõ rệt cho các bên liên quan trực tiếp đến hoạt động kinh doanh và vận hành hệ thống. 

\subsubsection{Khách hàng}

Ở cấp độ quản lý, hệ thống này đóng vai trò then chốt trong việc củng cố khả năng cạnh tranh của doanh nghiệp trên thị trường số. Việc cung cấp trải nghiệm AI tiên tiến giúp thu hút và giữ chân khách hàng, đồng thời tăng tỷ lệ chuyển đổi đơn hàng thông qua gợi ý cá nhân hóa. Về mặt kỹ thuật, hệ thống đảm bảo tính ổn định, bảo mật và khả năng mở rộng, giúp doanh nghiệp giảm thiểu rủi ro vận hành và dễ dàng phát triển dịch vụ trong tương lai mà không cần thay đổi kiến trúc cốt lõi.

\subsubsection{Đội ngũ vận hàng và quản trị viên}

Đối với đội ngũ nội bộ, hệ thống mới giúp tăng hiệu suất làm việc và tối ưu hóa quy trình quản lý. Lợi ích chính là việc tự động hóa khâu chăm sóc khách hàng, giúp đội ngũ hỗ trợ giảm tải khối lượng câu hỏi thường gặp. Quản trị viên có được một giao diện tập trung để quản lý sản phẩm, giá cả và khuyến mãi hiệu quả. Hơn nữa, hệ thống cung cấp báo cáo cơ bản và cơ chế cập nhật trạng thái đơn hàng, hỗ trợ công tác quản lý và ra quyết định.

\subsubsection{Chủ sở hữu doanh nghiệp}

Ở cấp độ quản lý, hệ thống này đóng vai trò then chốt trong việc củng cố khả năng cạnh tranh của doanh nghiệp trên thị trường số. Việc cung cấp trải nghiệm AI tiên tiến giúp thu hút và giữ chân khách hàng, đồng thời tăng tỷ lệ chuyển đổi đơn hàng thông qua gợi ý cá nhân hóa. Về mặt kỹ thuật, hệ thống đảm bảo tính ổn định, bảo mật và khả năng mở rộng, giúp doanh nghiệp giảm thiểu rủi ro vận hành và dễ dàng phát triển dịch vụ trong tương lai mà không cần thay đổi kiến trúc cốt lõi.

\subsubsection{Nhóm phát triển đề tài}

Dự án là nền tảng để nhóm phát triển ứng dụng và kiểm chứng các kỹ thuật hiện đại, đặc biệt là trong lĩnh vực Xử lý Ngôn ngữ Tự nhiên (NLP) và thiết kế kiến trúc vi dịch vụ cho thương mại điện tử. Lợi ích là kinh nghiệm thực tiễn trong việc tích hợp thành công công nghệ AI vào một hệ thống giao dịch phức tạp, tạo ra một sản phẩm minh chứng có giá trị nghiên cứu và ứng dụng cao.

Từ việc phân tích các bên liên quan, có thể thấy nhu cầu của từng nhóm được cụ thể hóa rõ ràng và gắn liền với vai trò mà họ đảm nhận trong hệ thống.Những nhu cầu trên chính là cơ sở để xác định các yêu cầu chức năng của hệ thống, bao gồm tính năng tìm kiếm, giỏ hàng, thanh toán mô phỏng, trợ lý ảo kịch bản và quản trị sản phẩm. Song song, các yêu cầu phi chức năng ở mức tối thiểu như sự ổn định cho bản demo, tốc độ phản hồi hợp lý và bảo mật thông tin người dùng cũng cần được đảm bảo để hệ thống có thể vận hành hiệu quả trong phạm vi nghiên cứu.